\begin{frame}[noframenumbering,plain]
  \setcounter{framenumber}{1}
  \maketitle
\end{frame}

\begin{frame}
  \frametitle{Положения, выносимые на защиту}
  \begin{enumerate}
    \item Математическая модель представления многопоточной программы
      в~виде частично упорядоченного множества операций с~памятью,
      построенная с~учётом модели памяти;
    \item Алгоритмы эффективного построения всех линейных порядков
      частично упорядоченного множества операций с~памятью с~учётом
      классов эквивалентности перестановок;
    \item Структура данных --- граф состояний (State~Graph) --- для
      компактного и~полного представления всех возможных результатов
      исполнения многопоточной программы;
    \item Алгоритмы верификации соответствия аппаратуры заявленной
      модели памяти как на основе линейных порядков, так и на основе
      анализа графа состояний.
  \end{enumerate}
\end{frame}
\note{
  Проговариваются вслух положения, выносимые на защиту.
}

\begin{frame}
  \frametitle{Содержание}
  \tableofcontents
\end{frame}
\note{
  Работа состоит из пяти глав.

  \medskip
  В первой главе анализируются современные подходы к~верификации
  подсистемы памяти с~учётом моделей памяти.

  Во второй главе рассматривается представление многопоточной программы
  в~виде частично упорядоченного множества операций с~памятью.

  Третья глава посвящена построению всех линейных порядков с~учётом
  зависимостей операций.

  В четвёртой главе предложена структура данных ``граф состояний''
  для компактного представления результатов исполнения.

  В пятой главе приведены результаты вычислительных экспериментов.
}
